\sectionwithtoclink{Introduction}

The Black--Scholes--Merton (BSM) model prices options under the assumption of constant volatility across strikes and maturities.
This implies a flat implied volatility surface which, in practice, is far from the case. Market-implied volatilities are both skewed and far from
constant across maturities, so any serious attempt to model option prices must therefore account for stochastic volatility.

The Heston model~\cite{heston1993} is one such model, extending BSM by letting the volatility follow a mean-reverting square-root (CIR) process
which is allowed to be correlated with the asset. One of the advantages of the Heston model is that it captures a lot of the shape of the volatility surface,
while still having a nice analytical expression for the characteristic function of the log asset price. The CF being well enough behaved turns out to be all one needs to price derivatives
due to Fourier methods, which we will use the beginning of this paper to show how.

This paper tries to develop as much of the pipeline as possible. After finding a pricing formula for calls when the CF has certain properties,
we spend some time deriving the original expression for the Heston CF and verify that it has the aforementioned properties. Moving on to implementation, we set
up an experiment to compare quadratures as well as the impact of introducing an analytical gradient~\cite{germano2017} on the speed of calibration. Finally,
we calibrate the model to S\&P 500 closing option prices under two weighting schemes and examine how the choice of objective function shapes the fitted implied volatility surface.



Notation Note: We take a lot of expectations in the beginning of the paper, usually of what turns out to be stochastic processes when taking these we know their values at time 0 making every expectation $\mathbb{E}_{(x_0,v_0)}$ 
instead of writing this every time we will simply define this as the default probabilitymeasure. Usually we like to note information up to a time as $\mathcal{F}_t$ in this paper i've choosen to drop the filters i.e work from the filter $\mathcal{F}_0$. One should 
note that since the heston model is a markov process everything still works as every expectation is essentially $\mathbb{E}_{(x_0,v_0)}$ and by the markov property 
$\mathbb{E}_\mu[f(X_{t+s},V_{t+s})|\mathcal{F}]=\mathbb{E}_{(x_t,v_t)}[f(X_s,V_s)]$ for $f$ non negative and measuerable, so this actually dosen't change anything simply set $T=T-t\quad x_0=x_t,v_0=v_t$




